\subsection{From Logistic Regression to Feedforward NN}
To fit a model to some data distribution, a simple and common approach is to use \textbf{Maximum Likelihood Estimation} (MLE).

\subsubsection{Maximum Likelihood Estimation}
In MLE, we want to estimate the parameters of an assumed probability distribution, given some observed data. For example, if we assume that the data is generated from a Gaussian distribution, we can estimate the mean and variance of that distribution using MLE.

More formally, given a dataset $\boldsymbol{X} = \{\boldsymbol{x}_1, \boldsymbol{x}_2, \ldots, \boldsymbol{x}_N\}$ and a known probability distribution $p_{\text{model}}(\boldsymbol{x}|\boldsymbol{\theta})$
we want to find the parameters $\boldsymbol{\theta}$ that better explain the data.
\begin{equation}
    \label{MLE}
    \boldsymbol{\theta}^* = \arg\max_{\boldsymbol{\theta}} p_{\text{model}}(\boldsymbol{X}|\boldsymbol{\theta})
\end{equation}
Assuming that the data points are independent and identically distributed (i.i.d.), we can express the likelihood as a product of individual probabilities:
\begin{equation*}
    p_{\text{model}}(\boldsymbol{X}|\boldsymbol{\theta}) = \prod_{i=1}^{N} p_{\text{model}}(\boldsymbol{x}_i|\boldsymbol{\theta})
\end{equation*}
To simplify the optimization process, we often take the logarithm of the likelihood function, which transforms the product into a sum:
\begin{equation*}
    \log p_{\text{model}}(\boldsymbol{X}|\boldsymbol{\theta}) = \sum_{i=1}^{N} \log p_{\text{model}}(\boldsymbol{x}_i|\boldsymbol{\theta})
\end{equation*}
We can show that this is equivalent to considering the expected value of the log-likelihood over the empirical (observed) data distribution:
\begin{equation*}
    \log p_{\text{model}}(\boldsymbol{X}|\boldsymbol{\theta}) = N \cdot \mathbb{E}_{\boldsymbol{x} \sim p_{\text{data}}}[\log p_{\text{model}}(\boldsymbol{x}|\boldsymbol{\theta})]
\end{equation*}
where $p_{\text{data}}$ is the empirical distribution of the observed data and $\mathbb{E}_{\boldsymbol{x} \sim p_{\text{data}}}$ denotes the expected value sampled from this distribution.

To find the optimal parameters $\boldsymbol{\theta}^*$, we usually differentiate the log-likelihood with respect to $\boldsymbol{\theta}$ and set it
to zero. We can then resolve the resulting equation to find the optimal parameters.

The equation \ref{MLE} can also be extended to the case of conditional probability distributions ($p_{\text{model}}(y|\boldsymbol{x}, \boldsymbol{\theta})$) in order to predict a target variable $y$ given an input variable $\boldsymbol{x}$:
\begin{equation*}
    \boldsymbol{\theta}^* = \arg\max_{\boldsymbol{\theta}} p_{\text{model}}(y|\boldsymbol{x}, \boldsymbol{\theta})
\end{equation*}

\subsubsection{Example: Logistic Regression}
One of the simplest examples of MLE is \textbf{Logistic Regression}, which is a linear model used for binary classification tasks. 
In Logistic Regression, we assume that the conditional probability distribution of the target variable $y$ given the input variable $\boldsymbol{x}$ follows a logistic function:
\begin{equation}
    \label{Logistic Regression}
    p_{\text{model}}(y=1|\boldsymbol{x}, \boldsymbol{\theta}) = \sigma(\boldsymbol{\theta}^T \boldsymbol{x}) = \frac{1}{1 + e^{-\boldsymbol{\theta}^T \boldsymbol{x}}}
\end{equation}
We can clearly see that the logistic function $\sigma(z)$ can be interpreted as a \textbf{neuron}:
\begin{itemize}
    \item The input to the neuron is the linear combination of the input features $\boldsymbol{x}$ and the parameters $\boldsymbol{\theta}$.
    \item The sigmoid activation function $\sigma(z)$ is a non-linear function that maps the input to a probability value between 0 and 1.
\end{itemize}
